% =====================================================================
%  Chapter 22 — Where it stands
% =====================================================================
\chapter{Where It Stands}
\label{ch:where-it-stands}

\begin{record}
\textbf{The one-paragraph answer.} The programme has not proved the Clay
Millennium Prize problem, and is not close to it. It claimed a proof in
\Sn{29} and withdrew that claim in \Sn{31} on eight grounds. Since then it has
built a substantial body of correct machinery around a single scalar quantity,
proved one hypothesis it had carried for nine sessions, proved seven things
impossible, and left one deep gap unmoved. Its best conceivable outcome, with
every remaining hypothesis discharged, is \emph{no Type-I singularity} ---
which \texttt{rem:s39-distance} states is ``strictly less than the Clay
problem'', and which the programme's own machinery cannot extend, because
without the Type-I rates there is no limit object at all.

\smallskip
That is the position. This chapter sets out what is proved, what is
conditional and on what, what is open, what is proved impossible, and what the
sharpest live target is --- which, as of \Sn{48}, is a question that did not
exist one session ago.
\end{record}

\section{What is proved}
\label{sec:proved}

Unconditional statements in the proof document, within their stated hypothesis
class. The list is not short, and that is worth saying: a programme can
accumulate a great deal of correct mathematics without approaching its goal.

{\footnotesize
\begin{longtable}{@{}p{0.30\textwidth} p{0.44\textwidth} l@{}}
\caption[What is proved]{Unconditional results, by session. ``Within their
stated hypothesis class'' does real work: most of the later entries assume the
Type-I bounds of line~5565, which are hypotheses about a putative singularity,
not theorems.}
\label{tab:proved}\\
\toprule
Object & Statement & Session \\
\midrule
\endfirsthead
\toprule
Object & Statement & Session \\
\midrule
\endhead
\bottomrule
\endfoot
\texttt{thm:routeC} & The \(\theta\)-bootstrap: any \(\delta_0>0\) gives
  smoothness. Still the terminal step of any closure, and never invoked
  since \Sn{31}, because no session produced a \(\delta_0\) & \Sn{17} \\
\texttt{thm:shear-claimA}, \texttt{thm:near-shear},
  \texttt{prop:lambda-zero} & Claim~A for shear and near-shear data;
  \(\lambda_{\max}=0\) for shear & \Sn{20}, \Sn{21} \\
\texttt{lem:grad-omega-split} & \(\abs{\nabla\omega}^2=\abs{\nabla m}^2
  +m^2\abs{\nabla\ehat}^2\). Load-bearing in \Sn{34} and \Sn{36} & \Sn{30} \\
Scale invariance of \(\sigma\) & The exponent \(3/2\) is the unique fixed
  point of the NS scaling group & \Sn{30} \\
\texttt{lem:direction-eq} & Constantin's magnitude and direction equations,
  re-derived; pressure-free & \Sn{32} \\
\texttt{prop:sigma-star-invariant}, \texttt{prop:sigma-controls} &
  \(\sstar\) exactly scale-invariant; \(\sstar\le16\sigma\) & \Sn{32} \\
\texttt{rem:marginality} & The direction equation is exactly critical at
  \(\rstar\), up to \(\Lfac\) & \Sn{32} \\
\texttt{prop:first-rung} & Angular Caccioppoli; routine items discharged in
  \Sn{36} & \Sn{33}, \Sn{36} \\
\texttt{thm:apriori-sigma} & \(\langle\sstar\rangle\le C(C_I)\Lfac/\nu\) near
  a Type-I singularity & \Sn{34} \\
\texttt{prop:crude-K} & The crude route misses closure by exactly one
  logarithm --- sharp & \Sn{34} \\
\texttt{prop:second-rung} & The critical quadratic is subordinate to
  \(\sstar\le\varepsilon_0\) & \Sn{35} \\
\texttt{lem:gram-negativity} & The drift--curvature couplings are
  good-signed & \Sn{36} \\
\texttt{lem:w2-cancel} & The \Sn{35} second good term cancels identically &
  \Sn{36} \\
\texttt{prop:localized-cf} & Localized Constantin--Fefferman at the shrinking
  scale (on \(\T^3\); \emph{not} on profiles --- Wall~8) & \Sn{37} \\
\texttt{prop:regime-A-impossible} & \(\limsup_{t\to T^-}\Omega=\infty\) at any
  genuine blow-up. Regime~A cannot occur & \Sn{38} \\
\texttt{lem:typeI-scaleinv-bounds} & The ancient Type-I class is scaling-closed
  with the same constant & \Sn{39} \\
\texttt{prop:transfer-audit}(a)--(c) & Magnitude, direction, split and Gram all
  hold on any profile, re-derived from scratch & \Sn{39} \\
\texttt{lem:alignment-local} & \(\sstar(U)=0\) forces \(\ehat\) constant
  locally & \Sn{39} \\
\texttt{lem:harmonic-reduction} & Aligned \(\Rightarrow\) \(\partial_3U=0\),
  \(U_h\) solves 2D NS. The dimensional reduction without which no Liouville
  theorem applies & \Sn{39}, \Sn{48} \\
\texttt{lem:sigma-semicontinuity} & \(\sstar(U)\le\liminf\sstar(u_{\lambda_j})\),
  and no reverse inequality & \Sn{39} \\
\texttt{prop:logtime-necessary} & The log-time averaged necessary condition
  (modulo attainment of the spatial supremum) & \Sn{39} \\
\texttt{prop:s41-core} & Six of the eight cutoff-derivative families are
  disposed of & \Sn{41} \\
\texttt{lem:s48-biotsavart} & \(\norm{v}_\infty\lesssim\norm{\omega}_\infty^{3/5}
  \norm{v}_{L^2}^{2/5}\), logarithm-free & \Sn{48} \\
\textbf{\texttt{thm:s48-R}} & \textbf{Hypothesis (R) is a theorem}, in a
  stronger form: \(\norm{\nabla^k\partial_s^lu_\lambda}_\infty\le
  C_{k,l}(c_I)\abs{s}^{-(k+1)/2-l}\), uniform in \(\lambda\) & \Sn{48} \\
\texttt{cor:s48-transfer-d} & \texttt{prop:transfer-audit}(d) first clause,
  unconditionally: the profile route never incurs \(\Lfac\) & \Sn{48} \\
\end{longtable}}

\section{What is conditional, and on what}
\label{sec:conditional}

\begin{center}
\footnotesize
\begin{tabular}{@{}p{0.30\textwidth} p{0.24\textwidth} p{0.36\textwidth}@{}}
\toprule
Statement & Conditional on & Note \\
\midrule
\texttt{thm:pincer-conditional}: every Type-I singularity maintains
  \(\sstar>\varepsilon_0\)
  & H1 (\texttt{conj:bridge}) and H2
  & H1's estimate layer is proved except \(\sstar\le\varepsilon_0\) and
    \(\Aann\); H2 in the generic regime \emph{is} \(\sstar\)-decay \\[3pt]
\texttt{prop:K-amended}: \(K\)-absorption with an absolute constant
  & \(\Aann\) (i.e.\ \texttt{target:band-absorption})
  & Reclassified from mechanical to open at \Sn{41} \\[3pt]
\texttt{prop:profile-extraction}, \texttt{lem:sigma-equality}
  & \textbf{(N)}
  & CKN gives the right bound only on \(Q_1(0)\), which touches the terminal
    slice \\[3pt]
\texttt{prop:transfer-audit}(e): \(\nu\langle\sstar\rangle\le C(c_I)\) on \(U\)
  & \textbf{saturation}
  & Was mod (R); (R) is now proved and this is what remains. See
    \S\ref{sec:saturation} \\[3pt]
\texttt{thm:aligned-horn}
  & \textbf{(A-up)}
  & The largest gap \Sn{39} created \\[3pt]
\texttt{prop:two-step-reduction}
  & \textbf{(N\('\))}
  & \\[3pt]
\texttt{thm:rigidity-dichotomy}
  & four inputs (was five)
  & (R) removed at \Sn{48} \\[3pt]
The whole \Sn{32}--\Sn{39} line
  & the Type-I rates of line~5565
  & \emph{Both} bounds, and the velocity one cannot be derived
    (\texttt{prop:s48-bridge-fails}) \\
\bottomrule
\end{tabular}
\end{center}

\section{What is open}
\label{sec:open}

\begin{center}
\footnotesize
\begin{tabular}{@{}p{0.26\textwidth} p{0.50\textwidth} l@{}}
\toprule
Item & What would have to be shown & Since \\
\midrule
\textbf{Wall 1}, the CZ logarithm
  & Any reformulation in which \(\nabla u\) on \(\T^3\) is not recovered by an
    \(L^\infty\) Calder\'on--Zygmund bound, or a proof that the logarithm can
    be beaten. No proof says it cannot & \Sn{31} (as D8) \\[3pt]
\texttt{target:disorder-depletion}
  & Persistent \(\sstar\ge\delta\) forces the log-time average
    \(\le1-c(\delta)\). Marginal: must beat \(1\) exactly & \Sn{39} \\[3pt]
\quad (T4a) & average-to-point for \(w\) at the vorticity maximum; weak
  Harnack is the natural mechanism & \Sn{39} \\[3pt]
\quad (T4b) & profile depletion of \(\alpha\); needs a far-field control
  \emph{other} than \texttt{prop:localized-cf}, which Wall~8 forbids on
  profiles & \Sn{39} \\[3pt]
\textbf{(A-up)} & local one-component alignment upgraded to global
  \(\omega\parallel e_0\) on \(\R^3\times(-\infty,0)\) & \Sn{39} \\[3pt]
\textbf{(N)}, \textbf{(N\('\))} & a Type-I-improved \(\varepsilon\)-regularity
  theorem, to push mass off the terminal slice & \Sn{39} \\[3pt]
\textbf{Saturation} & \(M_U(s)\gtrsim\abs{s}^{-1}\) for all \(s<0\), not only
  on a set; equivalently \(M(t)\gtrsim(T-t)^{-1}\) on \(\T^3\) & \Sn{48} \\[3pt]
\texttt{conj:K-refined}, \texttt{target:band-absorption},
  \texttt{conj:lc-s44}
  & each is a decorrelation hypothesis with \(C/\Lfac\) decay; (LC) covers
    both \(Y\)-halves & \Sn{34}, \Sn{41}, \Sn{44} \\[3pt]
The \(D\)-half, \eqref{eq:bandD}
  & a bare level non-concentration statement, covered by \emph{no} hypothesis
    in the family & \Sn{44} \\[3pt]
Type-II exclusion
  & outside all present methods; see \S\ref{sec:ceiling-vs-prize} & \Sn{35} \\
\bottomrule
\end{tabular}
\end{center}

\section{What is proved impossible}
\label{sec:impossible}

Each of these is \emph{qualified}, and the qualification is part of the
result. A programme that reported them unqualified would be overstating its
own negative results, which is the same failure mode as overstating the
positive ones.

\begin{center}
\footnotesize
\begin{tabular}{@{}p{0.29\textwidth} p{0.29\textwidth} p{0.32\textwidth}@{}}
\toprule
Result & Impossible \emph{as what} & What remains permitted \\
\midrule
Regime~A (bounded enstrophy near a singularity),
  \texttt{prop:regime-A-impossible}
  & Absolutely, for any genuine blow-up
  & Nothing; the route is closed. \(\sstar\)-decay is now necessary \\[3pt]
Nested level iteration, \texttt{prop:s41-nocontraction},
  \texttt{prop:s41-vacuous}, \texttt{prop:s41-coarea}
  & For the route, and for \emph{any} level selection
  & \texttt{target:band-absorption}, which is \OPEN{} \\[3pt]
\(c_K\Leftrightarrow\) band absorption,
  \texttt{prop:s44-a-fails}, \texttt{prop:s44-b-fails}
  & Over the a priori structure (A1)--(A3) only
  & An implication using further PDE structure. \textbf{Not} independence for
    Navier--Stokes solutions \\[3pt]
\(c_K\) numerics as evidence, Wall~6
  & As a test of the conjectures
  & Measurements whose answer differs between mechanisms in a
    \emph{regular} flow \\[3pt]
\texttt{prop:localized-cf} on profiles,
  \texttt{rem:s39-nontransfer}(i)
  & As written; the obstruction is a divergence
  & Some localized substitute on \(\R^3\); untried \\[3pt]
De-forcing a forced blow-up, Wall~9
  & As an import route
  & The \emph{geometry}, as unforced adversarial initial data \\[3pt]
The \(\omega\to u\) bridge,
  \texttt{prop:s48-bridge-fails}
  & Permanently; the exponent \(3/5\) is forced by scaling
  & Nothing. The velocity Type-I bound stays a hypothesis \\
\bottomrule
\end{tabular}
\end{center}

\section{The ceiling versus the Prize}
\label{sec:ceiling-vs-prize}

\begin{record}
Suppose everything in \S\ref{sec:open} were proved tomorrow: \textbf{(A-up)},
\texttt{target:disorder-depletion}, \textbf{(N)}, \textbf{(N\('\))},
saturation, and the band absorption target. The conclusion would be:
\emph{no Type-I singularity, at a non-degenerate point}.

\smallskip
That is not the Clay problem. The gap is not a missing lemma; it is the
construction. Every estimate from \Sn{33} to \Sn{39} uses the Type-I rates,
and the profile extraction uses them \emph{essentially} --- without
\(M(t)\lesssim(T-t)^{-1}\) the rescaled family has no uniform bound and there
is no limit object at all. In the document's words: ``A Type-II singularity
would rescale to nothing, and the entire reframe would be vacuous against
it.''
\end{record}

\begin{figure}[ht]
\centering
\resizebox{\textwidth}{!}{%
\begin{tikzpicture}[
  >={Latex[length=2mm]},
  bx/.style={draw, rounded corners=2pt, font=\scriptsize, align=center,
             inner sep=4pt, text width=38mm, minimum height=11mm},
  done/.style={bx, draw=provedgreen, very thick},
  op/.style={bx, draw=openblue, thick, fill=panelgrey},
  tii/.style={bx, draw=impossiblered, thick, dashed},
  ar/.style={->, thick}
]
\node[done] (mach) at (0,0)
  {\textbf{Proved machinery}\\ \Sn{32}--\Sn{48}\\
   \(\sstar\), the two rungs, the profile, (R)};
\node[op] (gaps) at (4.9,0)
  {\textbf{Open}\\ (A-up), (N), (N\('\)),\\ saturation,
   \texttt{target:disorder-}\\\texttt{depletion}, Wall~1};
\node[done] (typeI) at (9.8,0)
  {\textbf{Best case}\\ no Type-I singularity\\ at a non-degenerate point};
\node[tii] (typeII) at (9.8,-3.0)
  {\textbf{Type-II}\\ outside all present methods\\
   ``would rescale to nothing''};
\node[bx, draw=impossiblered, very thick, text width=38mm] (prize) at (14.4,-1.5)
  {\textbf{The Clay problem}\\ (A)/(B): global regularity,\\ unforced};

\draw[ar] (mach) -- (gaps);
\draw[ar] (gaps) -- (typeI);
\draw[ar, provedgreen] (typeI) -- (prize);
\draw[ar, impossiblered, dashed] (typeII) -- (prize);
\node[font=\scriptsize, text=impossiblered, align=center, text width=30mm]
  at (7.0,-3.5) {no route from here};

\node[font=\scriptsize\itshape, text=rulegrey, align=center, text width=104mm]
  at (6.4,-5.0)
  {Both arrows into the Prize are required. The programme has a plan for the
   upper one and no method at all for the lower one, and the reason is
   structural: the compactness that powers the whole construction is powered
   by the Type-I rates.};
\end{tikzpicture}}
\caption[The distance to the Prize]{Wall~3 in one picture. Closing every open
item in \S\ref{sec:open} delivers the upper arrow only.}
\label{fig:distance}
\end{figure}

\begin{obsv}[Wall 3 has an example now]
\label{obs:wall3-example}
Until recently Type-II was an abstract reservation. It is not any more. The
blow-up profile described in Chapter~\ref{ch:openai} has azimuthal component
\(u_\vartheta\sim\tau^{-(1/2+h)}\) with \(h>0\), which diverges strictly
faster than velocity Type-I --- a Type-II profile. The one explicitly
constructed singular geometry now in the literature for a
Navier--Stokes-type system sits exactly in this programme's blind spot. That
cuts against the programme and is recorded for that reason.
\end{obsv}

\section{The sharpest live target: saturation}
\label{sec:saturation}

\begin{record}
\textbf{The question.} Is the Type-I \emph{upper} bound matched by a lower
bound of the same order? On the profile: does \(M_U(s)\gtrsim\abs{s}^{-1}\)
hold for all \(s<0\), or only on a set? On the torus: does
\(M(t)\gtrsim(T-t)^{-1}\)?
\end{record}

Three properties make this the target the programme should attack next, and
all three are recorded in the \Sn{48} residue and scope remarks.

\textbf{It is the common requirement of both branches.} On the profile branch,
\texttt{prop:transfer-audit}(e) is \PROVEDMOD{saturation}: it uses (d) in the
form \(\norm{\nabla U}_\infty\lesssim M_U\), which (d) supplies only on the
saturated set, the two forms differing by \(\left(\abs{s}M_U(s)\right)^{-1}\).
On the torus branch, \texttt{obs:s48-wall1} gives a logarithm-free
\(\norm{\nabla u}_{L^\infty}\le C(c_I)(T-t)^{-1}\) --- which does not remove
Wall~1, because \S33--\S37 need \(\nabla u\) measured against \(M(t)\), and
the two agree \emph{only on the saturated set}. It is the same proviso, one
level apart.

\textbf{It is a statement about one scalar function of one variable.} Not a
recursion, not a multi-term budget, not a marginal comparison of averages: the
behaviour of \(M_U(s)\) against \(\abs{s}^{-1}\) on an ancient solution. Under
\textbf{(N)} one already has \(M_U(s)>0\) for every \(s<0\); what is needed is
the quantitative strengthening from positivity to saturation.

\textbf{It did not exist as a named question before \Sn{48}.} It is not a
target that has resisted attack; it is a target nobody has attacked. That is
not evidence it is easy. It is a reason to look.

The \Sn{48} scope remark puts the reframe in one sentence: the session
``replace[s] the question `is the profile route logarithm-free?' with the
sharper question of residue~(1): \emph{is the Type-I upper bound
saturated?}''

\subsection{The other live targets, in order}

\begin{enumerate}[leftmargin=2.2em]
\item \textbf{Saturation}, as above. Cheapest, best-defined, newest.
\item \textbf{(A-up)}, the global alignment upgrade. Three mechanisms are on
  record; the unique-continuation route reuses the programme's own \(K\)-family
  machinery and is judged mildly positive.
\item \textbf{(T4a)}, average-to-point control of \(w\) at the vorticity
  maximum, by weak Harnack. It plugs into machinery that exists.
\item \textbf{Wall 1}, on the torus branch. \OPEN{} and structural. The
  permitted direction is a reformulation in which \(\nabla u\) is not recovered
  by an \(L^\infty\) Calder\'on--Zygmund bound --- for instance a
  symmetry-reduced setting where it comes from lower-dimensional elliptic
  problems. That is untested and must be settled by a written estimate before
  anything is built on it.
\item \textbf{(T4b)} and \texttt{target:disorder-depletion} proper. Marginal,
  and \Sn{39} named marginality as the pattern in which this programme's
  errors have historically hidden.
\end{enumerate}

Two directions are explicitly \emph{not} live. Further \(c_K\)-family numerics
as evidence are forbidden by Wall~6. Any attempt to weaken the standing
hypotheses to a vorticity-only Type-I bound is forbidden permanently by
\texttt{prop:s48-bridge-fails}.

\section{A note on what this work should be called}
\label{sec:naming}

\begin{record}
If the claim discussed in Chapter~\ref{ch:openai} --- finite-time blow-up for
the \emph{forced} three-dimensional Navier--Stokes equations, establishing
Fefferman's alternative (C) --- is accepted by the community, then ``an
attempt on the Navier--Stokes Millennium Prize problem'' stops being an
accurate description of this programme.

\smallskip
The accurate description is: \emph{an attempt on global regularity for the
unforced equations}, Fefferman's alternatives (A) and (B).

\smallskip
The target does not move. Only the name does.
\end{record}

Three things follow, and none of them is a consolation.

First, (A) and (C) are \emph{not mutually exclusive}. Unforced global
regularity and forced breakdown can both be true, and most experts expect
exactly that. A proof of (C) wins the Prize and leaves (A) open as a
mathematical question --- the question this programme has actually been
attacking, in every chapter, without ever once introducing a forcing term.

Second, ``it is only (C) and (D)'' is not a refutation of anyone, and this
book has no interest in making it sound like one. Fefferman offered (C) and
(D) deliberately, ``to give reasonable leeway to solvers while retaining the
heart of the problem''. Taking the leeway is the stated rules. An earlier
draft of Chapter~\ref{ch:openai} treated the point as though it were itself
the argument; it is not, and the correct argument is the first one above,
which is stronger because it does not require anybody to be wrong.

Third, this book reports and does not adjudicate. A Lean formalisation is
strong evidence about a proof, but what it certifies is the statement
\emph{as formalised}; whether that statement is Fefferman's (C) is a question
for the community. Nothing in this programme depends on the outcome.

\section{What the founding thesis came to}
\label{sec:thesis-verdict}

Chapter~\ref{ch:origin} states it, and it should be restated here rather than
left in the first chapter where a reader will have forgotten it.

The analytical line that produced this programme's strongest results
\emph{does not use temperature}. From \Sn{18} onward the proof document works
with the exact incompressible Navier--Stokes equations at constant viscosity
and with the auxiliary scalar \(\theta\); the word ``temperature'' does not
appear in its mathematics. What survived the founding thesis is a method, not
a mechanism: find a scalar quantity that carries the dynamics, and use the
tools only scalars admit. The programme applied that method to \(T\), to
\(\theta\), to \(\abs{\omega}\), and to \(\sstar\), and only the last two are
still load-bearing.

The thesis was right about the shape of the tool and wrong about which scalar
to put in it. That is a real finding and it is worth having; it is not what
the programme set out to find.

\section{The final ledger}
\label{sec:final-ledger}

\begin{table}[ht]
\centering
\small
\begin{tabular}{@{}l r p{0.44\textwidth}@{}}
\toprule
Category & Count & Comment \\
\midrule
Sessions & \(47\) & numbered to \Sn{48}; no \Sn{19} \\
Proof document & \(11{,}309\) lines & \(152\) pages; first compiled at \Sn{47} \\
Logged experiments & \(124\) & across \(31\) database tables \\
Walls & \(10\) & 3 \WITHDRAWN{}, 4 \IMPOSSIBLE{} (all qualified), 3 \OPEN \\
Prize claims made & \(1\) & \Sn{29} \\
Prize claims withdrawn & \(1\) & \Sn{31}, on eight grounds \\
Hypotheses discharged & \(1\) & \textbf{(R)}, at \Sn{48} \\
Open items in the ledger & \(10\) & \S\ref{sec:open} \\
Distance to the Prize & --- & \emph{unchanged since \Sn{39}} \\
\bottomrule
\end{tabular}
\caption[The final ledger]{The programme in numbers. The last row is the one
that matters, and it is the one no session since \Sn{39} has altered ---
including \Sn{48}, which says so itself.}
\label{tab:final-ledger}
\end{table}

\begin{obsv}[The honest closing position]
\label{obs:closing}
The programme is further from the Clay problem than it believed itself to be
in \Sn{29}, and knows its position far more precisely than it did then. Those
two facts are the same fact. Every session since the audit has either added
correct machinery or removed a route, and removing routes has been the more
common outcome: of the ten walls, four carry proofs, and three of those four
were produced by sessions commissioned to do the opposite.

\smallskip
What the programme has, at the end, is a scalar quantity with an exact scale
invariance, a complete estimate layer around it modulo two named hypotheses, a
limit-profile framework on which one horn of a dichotomy is closed but for a
single propagation hypothesis, a theorem that removes a hypothesis carried for
nine sessions, and a precise, short, unattacked question --- saturation ---
that both remaining branches turn out to need.

\smallskip
What it does not have is any route to Type-II, any proof that the
Calder\'on--Zygmund logarithm can be beaten on the torus, or any claim on the
Prize. The ceiling of \texttt{rem:s39-distance} stands verbatim, nine sessions
after it was written, and it will stand until somebody finds a method that
does not begin by rescaling.
\end{obsv}
